\documentclass[12pt]{article}
\usepackage[a3paper, margin=1in]{geometry}
\usepackage{amsmath, amssymb, amsthm, ulem}

\begin{document}

\section*{Domácí úkol 7.2}

Vypracoval: Daniel \textit{``Randál"} Ransdorf \hfill Podpis: \rule{4cm}{0.4pt}

\begin{enumerate}

  \setcounter{enumi}{1}
  
  \item Chceme zjistit, zda se dá nulový vektor netrivialně zapsat jako lineární kombinace
    vektorů z posloupnosti. Zkontrolujme tedy existenci nenulové trojice koeficientů $a,b,c\in\mathbb{R}$.

    \begin{align*}
      a(4\vec{u}+3\vec{v}+2\vec{w}+\vec{z}) + b(\vec{u}-\vec{v}-\vec{w}+\vec{z}) + c(\vec{u}+\vec{v}+\vec{w}+\vec{z}) &= \vec{0} \\
      \left(\begin{array}{ccc}
        4\vec{u}+3\vec{v}+2\vec{w}+\vec{z} & \vec{u}-\vec{v}-\vec{w}+\vec{z} & \vec{u}+\vec{v}+\vec{w}+\vec{z}
      \end{array}\right) \cdot \left(\begin{array}{c}
        a\\b\\c
      \end{array}\right) &= \vec{0} \\
      \overset{\text{řádkový pohled}}{\Longrightarrow}
      \left(\begin{array}{cccc}
        \vec{u} & \vec{v} & \vec{w} & \vec{z}
      \end{array}\right) \cdot \left(\begin{array}{ccc}
        4 & 1 & 1 \\ 3 & -1 & 1 \\ 2 & -1 & 1 \\ 1 & 1 & 1
      \end{array}\right) \cdot \left(\begin{array}{c}
        a\\b\\c
      \end{array}\right) &= \vec{0} \\
      \text{Položme } \vec{\alpha} = \left(\begin{array}{ccc}
        4 & 1 & 1 \\ 3 & -1 & 1 \\ 2 & -1 & 1 \\ 1 & 1 & 1
      \end{array}\right) \cdot \left(\begin{array}{c}
        a\\b\\c
      \end{array}\right)
    \end{align*}
    Z lineární nezávislosti víme, že nulový vektor z posloupnosti ($\vec{u},\vec{v},\vec{w},\vec{z}$) nemůžeme vyrobit
    lineární kombinací jinak, než součtem 0-násobků daných vektorů. Tedy aritmetický vektor $\vec{\alpha} = \vec{0}$ je jediné řešení rovnice:
    \begin{align*}
      \left(\begin{array}{cccc}
        \vec{u} & \vec{v} & \vec{w} & \vec{z}
      \end{array}\right) \cdot \vec{\alpha} = \vec{0} \\
    \end{align*}
    Díky tomu můžeme problém redukovat na rovnici $\vec{\alpha} = \vec{0}$:
    \begin{align*}
      &\left(\begin{array}{ccc}
        4 & 1 & 1 \\ 3 & -1 & 1 \\ 2 & -1 & 1 \\ 1 & 1 & 1
      \end{array}\right) \cdot \left(\begin{array}{c}
        a\\b\\c
      \end{array}\right) = \vec{0} \\
      &\left(\begin{array}{ccc|c}
        4 & 1 & 1 & 0 \\ 3 & -1 & 1 & 0 \\ 2 & -1 & 1 & 0 \\ 1 & 1 & 1 & 0
      \end{array}\right)
      \overset{\text{Gauss--Jordan eliminace}}{\sim}
      \left(\begin{array}{ccc|c}
        1&0&0&0 \\ 0&1&0&0 \\ 0&0&1&0 \\ 0&0&0&0
      \end{array}\right) \\
      \Longrightarrow
      &\left(\begin{array}{c}
        a\\b\\c
      \end{array}\right) = \left(\begin{array}{c}
        0\\0\\0
      \end{array}\right)
    \end{align*}

    Vyřešili jsme jediné možné koeficienty $a,b,c$ a všechny jsou rovny $0$. To znamená,
    že \textbf{jediný} způsob, jak ze zadané posloupnosti vytvořit lineární kombinací nulový vektor je triviálně nastavit koeficienty rovny $0$.
    Tím pádem je posloupnost \underline{lineárně nezávislá} ve \textbf{V}.
    

\end{enumerate}

\end{document}
